% !TEX program = pdflatex
%
% example.tex: a paper with many comments, as it may look before it is
% uploaded to arXiv. Switch the steps on and off to see what they remove.
%
\documentclass{article}
\usepackage{amsmath,amssymb}
\usepackage{listings}
\usepackage{hyperref}

% TODO: pick a better title
\newcommand{\R}{\mathbb{R}}% the real numbers
\newcommand{\norm}[1]{%
  \left\lVert #1 \right\rVert% no space here
}

\title{Removing Comments}  % working title
\author{A. Author}

\begin{document}
\maketitle

\section{Introduction}
% Reviewer 2 asked for more motivation.
% We should cite Knuth here.
Comments are useful while writing, % but not in the final version
but they end up in the source that you upload.	% e.g., to arXiv
About 20\% of this file are comments.

%\section{Old section}
%This section was removed after the first review.


\section{Results}
The norm is
\[
  \norm{x} = \sqrt{x_1^2 + x_2^2} % Euclidean
  % \norm{x}_1 = |x_1| + |x_2|
  % was the first try
\]
for $x \in \R^2$. %TODO: generalize

\begin{tabular}{ll}
  Step & Result \\ % header
  1 & 20\% \\
\end{tabular}

Mr.\ % a control space keeps its space
Smith writes \verb|%| in the code, see \url{https://example.org/a%20b}.

\begin{lstlisting}[language=Python]
x = 7 % 3  # the modulo operator, not a comment
\end{lstlisting}

\end{document}
